\documentclass[../../../include/open-logic-section]{subfiles}

\begin{document}

\olfileid{sth}{choice}{woproblem}
\olsection{సుక్రమపరచడం సమస్య}

సుక్రమపరచడాన్ని అంగీకరిస్తామా లేదా అన్నదానిపై చాలా విషయాలు ఆధారపడతాయని స్పష్టం. అయితే ఈ సూత్రంపై చర్చ ఎక్కువగా దానికి సమానమైన ఎంపిక స్వయంసిద్ధంపై కేంద్రీకృతమైంది. కాబట్టి ఇప్పుడు దానిని పరిశీలించి, సమానత్వాన్ని నిరూపిద్దాం.

\citeyear{Cantor1883}లో కాంటర్ సుక్రమపరచడం స్వయంసిద్ధానికి మద్దతు తెలుపుతూ, దాన్ని ``నాకు ప్రాథమికంగా, విస్తృత పరిణామాలు కలిగినదిగా, ముఖ్యంగా సార్వత్రిక చెల్లుబాటు వల్ల విశేషమైనదిగా కనిపించే ఆలోచనా నియమం'' అన్నారు (\citeauthor{Potter2004}లో ఉటంకించబడింది, \citeyear[p.~243]{Potter2004}). కానీ చివరికి ఆ ``స్వయంసిద్ధం'' నిరూపణ కావాలని కాంటర్ నమ్మారు. గణిత సమాజమూ అలాగే భావించింది.

ఈ సమస్యను \citeyear{Zermelo1904}లో జెర్మెలో ``పరిష్కరించారు''. ఆయన పరిష్కారాన్ని వివరించడానికి కొన్ని నిర్వచనాలు కావాలి.

\begin{defn}
ప్రతి $x \in \dom{f}$కు $f(x) \in x$ అయితే, అప్పుడే $f$ అనే ప్రమేయాన్ని \emph{ఎంపిక ప్రమేయం} అంటాం. $f$ ఎంపిక ప్రమేయమై, $\dom{f} = A \setminus \{\emptyset\}$ అయితే, అప్పుడే $f$ను $A$ కోసం \emph{ఎంపిక ప్రమేయం} అంటాం.
\end{defn}

సహజంగా, ప్రతి ఖాళీ కాని సమితి $x \in A$ నుంచి $A$ కోసం ఎంపిక ప్రమేయం ఒక నిర్దిష్ట మూలకం $f(x)$ను అదే $x$ నుంచి \emph{ఎంచుకుంటుంది}. అప్పుడు ఎంపిక స్వయంసిద్ధం ఇది:

\begin{axiom}[ఎంపిక]
	ప్రతి సమితికీ ఎంపిక ప్రమేయం ఉంది.
\end{axiom}

ఎంపిక వల్ల సుక్రమపరచడం వస్తుందని, అలాగే తిరుగుదిశలోనూ వస్తుందని జెర్మెలో చూపారు:

\begin{thm}[$\ZF$లో]\ollabel{thmwochoice}
సుక్రమపరచడం, ఎంపిక సమానార్థకాలు.
\end{thm}

\begin{proof}
\emph{ఎడమ నుంచి కుడికి.} $A$ సమితుల సమితి అనుకుందాం. సమ్మేళన స్వయంసిద్ధం ప్రకారం $\bigcup A$ ఉంది; సుక్రమపరచడం ప్రకారం $\bigcup A$ను సుక్రమపరిచే $<$ సంబంధం ఉంది. ఇప్పుడు $f(x) = \text{$<$-అత్యల్ప మూలకం }x$ అని తీసుకుందాం. ఇది $A$ కోసం ఎంపిక ప్రమేయం.

\emph{కుడి నుంచి ఎడమకు.} $A$ను స్థిరపరచుకుందాం. అది ఖాళీ అయితే సుక్రమం స్పష్టం; కనుక ఇక ఖాళీ కానిదని అనుకుందాం. ఎంపిక ప్రకారం $\Pow{A} \setminus \{\emptyset\}$ కోసం $f$ అనే ఎంపిక ప్రమేయం ఉంది. అతిపరిమిత పునరావృత్తి ద్వారా ఈ విధంగా నిర్వచిద్దాం:
\sourcecorrection{OLTESTCHOICEWO-001}{మూల పునరావృత్తిలో $g(0)=f(A)$ అని ఉంది; $A$ ఖాళీ అయితే $f(A)$ నిర్వచితం కాదు. ఖాళీ సమితికి సుక్రమం స్పష్టమని విడిగా చెప్పి, పునరావృత్తికి ఖాళీ కాని $A$ను మాత్రమే తీసుకున్నాం.}
\begin{align*}
	g(0) &= f(A)\\
	g(\alpha) &= 
		\begin{cases}
			\text{ఆపు!{}} &\text{ఒకవేళ }A = \funimage{g}{\alpha}\\
			f(A \setminus \funimage{g}{\alpha}) & \text{లేకపోతే}\\	
		\end{cases}
\end{align*}
``ఆపు!'' అనేది ఎక్కువ మాటలతో చెప్పాల్సిన నిర్వచనానికి సంక్షిప్త సూచన మాత్రమే. అంటే మొదటిసారి $A =
\funimage{g}{\alpha}$ అయినప్పుడు, ఆ దశ నుంచి అన్ని $\delta \geq \alpha$లకు $g(\delta) = A$ అనే ఆపు-గుర్తును ఉంచి, తరువాతి వాదనలో ఆ దశకు ముందున్న ఎంపికల భాగాన్నే తీసుకుంటాం. మొదట మనకు ఇది ఒక \emph{పదాన్ని} మాత్రమే నిర్వచిస్తుందని చెప్పగలం (\olref[sth][spine][recursion]{transrecursionschema} తరువాతి వ్యాఖ్యలు చూడండి); కానీ ఆపే దశలోని పరిమితి నిజంగా ప్రమేయం అవుతుందని చూపుతాం.
\sourcecorrection{OLTESTCHOICEWO-002}{మూలం ఆపే దశ తరువాతి గుర్తును అన్ని $\delta\leq\alpha$కు పెట్టింది; అలా చేస్తే ముందుగా ఎంచుకున్న విలువలన్నీ చెరిగి ఇంజెక్టివ్ వాదన విఫలమవుతుంది. గుర్తును $\delta\geq\alpha$కు మాత్రమే ఉంచి, ఆపే ముందరి భాగమే అవసరమైన ప్రమేయమని స్పష్టం చేశాం.}

$f$ ఎంపిక ప్రమేయం కాబట్టి, ఆపే ముందు నిర్వచితమైన ప్రతి $\alpha$కు $g(\alpha) =
f(A \setminus \funimage{g}{\alpha}) \in A \setminus
\funimage{g}{\alpha}$; అంటే $g(\alpha) \notin \funimage{g}{\alpha}$. కాబట్టి $g(\alpha) = g(\beta)$ అయితే $g(\beta) \notin
\funimage{g}{\alpha}$; అంటే $\beta \notin \alpha$. అలాగే $\alpha \notin \beta$. త్రివిభాగ ధర్మం వల్ల $\alpha = \beta$. కనుక ఆపే ముందరి $g$ \tetoken{ఇంజెక్టివ్}{!!{injective}}.

తరువాత, నిజంగానే ఆపే దశ ఉందని గమనించండి: $A = g[\alpha]$ అయ్యే ఏదో అత్యల్ప క్రమసంఖ్య $\alpha$ ఉంది. కాదనుకుంటే, $g$ \tetoken{ఇంజెక్టివ్}{!!{injective}} కాబట్టి ప్రతి క్రమసంఖ్య $\alpha$కూ $\cardless{\alpha}{\Pow{A} \setminus \{\emptyset\}}$ వస్తుంది; ఇది \olref[hartogs]{HartogsLemma}కు విరుద్ధం. అందువల్ల ఆపే ముందు భాగానికి $\ran{g} = A$ కూడా.

ఈ విషయాలను కలిపితే, ఆపే ముందు భాగమైన $g$ ఏదో క్రమసంఖ్య నుంచి $A$కు \tetoken{ద్వైజెక్షన్}{!!a{bijection}}. ఇప్పుడు $g$ను వాడి $A$ను సుక్రమపరచవచ్చు.
\end{proof}

కాబట్టి సుక్రమపరచడం, ఎంపిక రెండూ కలిసే నిలుస్తాయి లేదా పడిపోతాయి. కానీ ప్రశ్న మిగిలే ఉంది: అవి నిలుస్తాయా, పడిపోతాయా?

\end{document}
